%%%%%%%%%%%%%%%%%%%%%%%%%%%%%%%%%%%%%%%%%%%%%%%%%%%%%%%%%%%%%%%%%%%%%%%%%%%%%%%%%%%%%%%%%%%%%%%%%%%%%%%%%
%%%%%%%%%%%%%%%%%%%%%%%%%%%%%%%%%%%%%%%%%%%%%%%%%%%%%%%%%%%%%%%%%%%%%%%%%%%%%%%%%%%%%%%%%%%%%%%%%%%%%%%%%

%% POJĘCIA
\newglossaryentry{kontener}{
    name={Kontener},
    description={Lekka, przenośna i skalowalna jednostka, która zawiera oprogramowanie oraz jego zależności.}
}

\newglossaryentry{kubernetes}
{
    name=Kubernetes,
    description={Otwartoźródłowa platforma do zarządzania, automatyzacji i skalowania aplikacji kontenerowych}
}


\newglossaryentry{controlplane}{
    name={Control Plane},
    description={Warstwa sterująca klastra Kubernetes, odpowiedzialna za zarządzanie stanem klastra, planowanie zadań i wykrywanie oraz reagowanie na zdarzenia w klastrze}
}

\newglossaryentry{workernode}{
    name={Worker Node},
    description={Węzeł roboczy w klastrze Kubernetes; maszyna fizyczna lub wirtualna, na której uruchamiane są Pody z aplikacjami użytkownika}
}

\newglossaryentry{pod}{
    name={Pod},
    description={Najmniejsza i najprostsza jednostka w modelu obiektowym Kubernetesa, którą można utworzyć lub wdrożyć. Pod reprezentuje pojedynczą instancję działającego procesu w klastrze}
}

\newglossaryentry{deployment}{
    name={Deployment},
    description={Obiekt API Kubernetesa, który zarządza replikowanymi aplikacjami, umożliwiając deklaratywne aktualizacje Podów i ReplicaSetów}
}

\newglossaryentry{service}{
    name={Service},
    description={Abstrakcja definiująca logiczny zbiór Podów oraz politykę dostępu do nich (często nazywaną mikroserwisem)}
}

\newglossaryentry{operator}{
    name={Operator},
    description={Metoda pakowania, wdrażania i zarządzania aplikacją Kubernetes, wykorzystująca niestandardowe zasoby (CRD) do zautomatyzowania zadań administracyjnych}
}

\newglossaryentry{metricsserver}{
    name={Metrics Server},
    description={Komponent klastra Kubernetes zbierający metryki zużycia zasobów (CPU, RAM) z poszczególnych węzłów i udostępniający je poprzez Metrics API}
}

\newglossaryentry{admissioncontroller}{
    name={Admission Controller},
    description={Fragment kodu, który przechwytuje żądania do serwera API Kubernetes przed utrwaleniem obiektu, ale po uwierzytelnieniu i autoryzacji. Może walidować (Validating) lub modyfikować (Mutating) obiekty}
}

\newglossaryentry{scaler}{
    name={Scaler},
    description={W terminologii KEDA: komponent definiujący logikę połączenia z konkretnym źródłem zdarzeń (np. Kafka, Prometheus) i metodę pobierania z niego metryki decydującej o skalowaniu}
}

\newglossaryentry{scaledobject}{
    name={ScaledObject},
    description={Niestandardowy zasób (CRD) wprowadzany przez KEDA, służący do konfiguracji automatycznego skalowania dla konkretnego wdrożenia (Deployment) w oparciu o zdefiniowane triggery}
}

\newglossaryentry{autonomicmanager}{
    name={Autonomic Manager},
    description={Komponent oprogramowania, który implementuje pętlę sterowania (MAPE-K) w celu zarządzania zasobem lub systemem bez interwencji człowieka}
}

\newglossaryentry{managedelement}{
    name={Managed Element},
    description={Element systemu (sprzęt lub oprogramowanie), który jest kontrolowany przez menedżera autonomicznego. W Kubernetesie jest to zazwyczaj Pod, Deployment lub Node}
}

\newglossaryentry{sensors}{
    name={Sensory},
    description={Interfejsy umożliwiające menedżerowi pobieranie informacji o stanie elementu zarządzanego (np. eksportery metryk w Prometheusie)}
}

\newglossaryentry{effectors}{
    name={Efektory},
    description={Interfejsy umożliwiające menedżerowi wprowadzanie zmian w elemencie zarządzanym (np. API Kubernetesa pozwalające na zmianę konfiguracji)}
}

\newglossaryentry{autonomiccomputing}{
    name={Autonomic Computing},
    description={Paradygmat systemów komputerowych zdolnych do samodzielnego zarządzania (samokonfiguracji, samoleczenia, samooptymalizacji i samoobrony) przy minimalnej interwencji człowieka}
}

% (Opcjonalnie, jeśli używasz w innym miejscu)
\newglossaryentry{feedbackloop}{
    name={Pętla sprzężenia zwrotnego},
    description={Mechanizm sterowania, w którym wyjście systemu jest mierzone i porównywane z wartością zadaną w celu wygenerowania sygnału sterującego minimalizującego błąd}
}


% W sekcji Glosariusza (jeśli używasz)
\newglossaryentry{thrashing}
{
    name=Thrashing,
    description={Zjawisko w systemach komputerowych, w którym zasoby są marnowane na ciągłe przełączanie kontekstu lub rekonfigurację (np. ciągłe skalowanie w górę i w dół) zamiast na wykonywanie użytecznej pracy}
}

\newglossaryentry{flash_crowd}
{
    name=Flash Crowd,
    description={Nagły, gwałtowny i często nieprzewidywalny wzrost ruchu sieciowego skierowanego do usługi, często prowadzący do przeciążenia zasobów}
}

\newglossaryentry{overhead}
{
    name=Narzut (Overhead),
    description={Dodatkowe zasoby obliczeniowe lub czasowe wymagane do wykonywania zadań zarządzania (np. monitorowania, podejmowania decyzji), które nie wynikają bezpośrednio z logiki biznesowej aplikacji}
}

% Glosariusz
\newglossaryentry{proxmox}
{
    name=Proxmox VE,
    description={Otwartoźródłowa platforma do zarządzania wirtualizacją serwerów, integrująca hiperwizor KVM i kontenery LXC}
}

\newglossaryentry{terraform}
{
    name=Terraform,
    description={Narzędzie typu open-source do tworzenia, zmieniania i wersjonowania infrastruktury w sposób bezpieczny i efektywny (IaC)}
}

\newglossaryentry{ansible}
{
    name=Ansible,
    description={Narzędzie do automatyzacji oprogramowania, zarządzania konfiguracją i wdrażania aplikacji}
}

\newglossaryentry{fluxcd}
{
    name=FluxCD,
    description={Narzędzie realizujące podejście GitOps dla platformy Kubernetes, zapewniające ciągłą synchronizację stanu klastra z repozytorium Git}
}

\newglossaryentry{k6}
{
    name=k6,
    description={Nowoczesne narzędzie do testowania obciążeniowego typu open-source, umożliwiające pisanie skryptów testowych w języku JavaScript}
}

\newglossaryentry{ubuntu}
{
    name=Ubuntu,
    description={Popularna dystrybucja systemu Linux oparta na Debianie, szeroko stosowana w środowiskach chmurowych}
}

\newglossaryentry{k3s}
{
    name=k3s,
    description={Lekka, certyfikowana dystrybucja Kubernetesa zaprojektowana z myślą o środowiskach IoT i Edge Computing, charakteryzująca się niskim zapotrzebowaniem na zasoby}
}

\newglossaryentry{flannel}
{
    name=Flannel,
    description={Prosta i lekka nakładka sieciowa (network overlay) dla Kubernetesa, realizująca model CNI}
}

\newglossaryentry{gitops}
{
    name=GitOps,
    description={Model operacyjny dla aplikacji cloud-native, w którym repozytorium Git stanowi "jedyne źródło prawdy" dla infrastruktury i aplikacji}
}

%%%%%%%%%%%%%%%%%%%%%%%%%%%%%%%%%%%%%%%%%%%%%%%%%%%%%%%%%%%%%%%%%%%%%%%%%%%%%%%%%%%%%%%%%%%%%%%%%%%%%%%%%
%%%%%%%%%%%%%%%%%%%%%%%%%%%%%%%%%%%%%%%%%%%%%%%%%%%%%%%%%%%%%%%%%%%%%%%%%%%%%%%%%%%%%%%%%%%%%%%%%%%%%%%%%


%% AKRONIMY 

\newacronym{vpa}{VPA}{Vertical Pod Autoscaler}
\newacronym{hpa}{HPA}{Horizontal Pod Autoscaler}
\newacronym{keda}{KEDA}{Kubernetes Event-driven Autoscaling}
\newacronym{crd}{CRD}{Custom Resource Definition}
\newacronym{api}{API}{Application Programming Interface}
\newacronym{rest}{REST}{Representational State Transfer}
\newacronym{cni}{CNI}{Container Network Interface}
\newacronym{cicd}{CI/CD}{Continuous Integration / Continuous Delivery}
\newacronym{mape}{MAPE}{Monitor, Analyze, Plan and Execute}
\newacronym{focale}{FOCALE}{Foundation - Observe - Compare - Act - Learn - rEason}
\newacronym{ooda}{OODA}{Observe-Orient-Decide-Act}
\newacronym{cladra}{CLADRA}{Closed Loop Anomaly Detection and Resolution Automation}
\newacronym{sla}{SLA}{Service Level Agreement}
\newacronym{k8s}{K8s}{Kubernetes}
\newacronym{mapek}{MAPE-K}{Monitor-Analyze-Plan-Execute-Knowledge}
\newacronym{iac}{IaC}{Infrastructure as Code}
\newacronym{qos}{QoS}{Quality of Service}

\newacronym{vm}{VM}{Virtual Machine}
\newacronym{cncf}{CNCF}{Cloud Native Computing Foundation}

